\section{Síntesis y ampliación}

\subsection{Puntos clave}

\begin{enumerate}
    \item La interacción activa es la base de la inteligencia: un LLM de solo texto es un ``gatito pasivo''
    \item Foundation Agent (Voyager, Eureka) muestra la exploración de mundo abierto impulsada por LLM y el diseño automático de reward
    \item Project GR00T busca construir un modelo base general entre distintas morfologías de robots
    \item La escasez de datos y el sim-to-real gap son los retos centrales
    \item La IA con cuerpo es un componente necesario de la AGI
\end{enumerate}

\subsection{Summary table}

\begin{table}[H]
\centering
\begin{tabular}{p{0.24\textwidth} p{0.2\textwidth} p{0.28\textwidth} p{0.2\textwidth}}
\toprule
Tema & Promesa & Reto & Producto clave \\
\midrule
Interacción activa con cuerpo & Ciclo cerrado de percepción-lenguaje-acción a la medida del robot & Escasez de datos + sim-to-real gap & Arquitectura de Voyager/Eureka, baseline de GR00T \\
Datos y entrenamiento & LLM-generated tasks + replay hygiene & Mantener la diversidad, distribución natural & Isaac Sim pipelines, generación de reward con LLM \\
Operación en producción & Trace + metrics + alerts & Cambio de herramientas, calibration drift & Incident playback, checklists \\
Gobernanza & Safety board + audit trail & Autonomous failure modes & Usage policy + override switch \\
\bottomrule
\end{tabular}
\caption{Principales dimensiones de aplicación práctica de la Lección 04}
\end{table}

\subsection{Lecturas adicionales}

\begin{itemize}
    \item Wang et al., ``Voyager: An Open-Ended Embodied Agent with Large Language Models,'' 2023.
    \item Ma et al., ``Eureka: Human-Level Reward Design via Coding Large Language Models,'' ICLR 2024.
    \item Fan et al., ``MineDojo: Building Open-Ended Embodied Agents with Internet-Scale Knowledge,'' NeurIPS 2022.
\end{itemize}


\end{document}
